\ticket{20}{Уровни языка и анализ текста}

\begin{examplan}
    \item Перечислить уровни естественного языка.
    \item Описать этапы многоуровневого анализа текста.
    \item Сравнить инженерный подход и подход на машинном обучении.
    \item Перечислить лингвистические ресурсы и их назначение.
\end{examplan}

\subsection*{Короткая версия для письменного ответа}

Обработка естественного языка обычно строится как многоуровневый анализ: от символов и слов к морфологии, синтаксису, семантике, дискурсу и прагматике. Каждый уровень добавляет информацию, которая может использоваться следующими модулями.

\subsection*{Уровни естественного языка}

\begin{itemize}
    \item \textbf{Фонетика и фонология}: звуки речи и их функции; важны для распознавания и синтеза речи.
    \item \textbf{Морфология}: структура слов, словоформы, леммы, грамматические признаки.
    \item \textbf{Лексика и лексическая семантика}: значения слов, синонимия, омонимия, многозначность.
    \item \textbf{Синтаксис}: структура словосочетаний и предложений.
    \item \textbf{Семантика}: смысл фраз, предложений и текста.
    \item \textbf{Прагматика}: влияние контекста, целей и ситуации общения.
    \item \textbf{Дискурс}: связность текста выше уровня предложения.
\end{itemize}

\subsection*{Модули анализа текста}

Типичный NLP-конвейер:
\begin{enumerate}
    \item сегментация текста на предложения;
    \item токенизация;
    \item нормализация, лемматизация или стемминг;
    \item POS-tagging и морфологический анализ;
    \item синтаксический анализ: зависимости или составляющие;
    \item распознавание именованных сущностей;
    \item разрешение многозначности;
    \item извлечение отношений и семантических ролей;
    \item кореференция, дискурс-анализ, прагматическая интерпретация.
\end{enumerate}

На практике порядок модулей зависит от задачи, а в нейросетевых системах несколько уровней могут обрабатываться одной моделью.

\subsection*{Инженерный подход}

\textbf{Подход на правилах} использует вручную заданные грамматики, регулярные выражения, словари и шаблоны.

Плюсы:
\begin{itemize}
    \item интерпретируемость;
    \item высокая точность в узких областях;
    \item не требует больших размеченных данных.
\end{itemize}

Минусы:
\begin{itemize}
    \item трудоемкость разработки;
    \item хрупкость к вариативности языка;
    \item плохая переносимость на новые домены.
\end{itemize}

\subsection*{Подход на машинном обучении}

ML-подход обучает модели на корпусах. Он может использовать обучение с учителем, без учителя, полуобучение и глубокие нейросетевые модели.

Плюсы:
\begin{itemize}
    \item лучше масштабируется на большие данные;
    \item улавливает сложные статистические закономерности;
    \item часто дает более высокое качество на стандартных задачах.
\end{itemize}

Минусы:
\begin{itemize}
    \item нужны данные и разметка;
    \item сложные модели менее интерпретируемы;
    \item качество сильно зависит от корпуса и домена.
\end{itemize}

\subsection*{Лингвистические ресурсы}

\begin{description}
    \item[Корпусы] большие коллекции текстов: неразмеченные, POS-размеченные, синтаксические treebanks, семантические, NER-корпусы, параллельные корпусы.
    \item[Лексические ресурсы] словари, тезаурусы, WordNet-подобные ресурсы, онтологии, списки именованных сущностей.
    \item[Грамматики и правила] формальные грамматики, морфологические правила, шаблоны извлечения.
\end{description}

Ресурсы нужны для обучения и оценки моделей, лингвистического анализа, информационного поиска, машинного перевода и извлечения знаний.

\begin{important}
    \item NLP анализирует текст на нескольких уровнях: морфология, синтаксис, семантика, дискурс.
    \item Конвейер обработки текста обычно начинается с токенизации и нормализации.
    \item Правила понятны и точны в узких задачах, но плохо масштабируются.
    \item ML-модели требуют данных, но лучше обобщают при хорошем обучении.
    \item Корпусы, словари, тезаурусы, онтологии и грамматики являются ключевыми лингвистическими ресурсами.
\end{important}

